\section{从多渠道到单 Agent：渠道坍缩与体验统一}

传统客服按渠道分工：电话、在线聊天、邮件、工单系统分别由不同团队与流程负责。讲者给出的方向是“single agent world”：一个可配置 Agent 在不同入口统一接待客户，再根据上下文、权限和任务类型调用不同动作。

这一策略不是简单“把多个入口接到同一个大模型”。真正难点是统一状态机和策略层：当用户从 chat 转到 voice，或者从 voice 回到工单，系统是否能保持任务上下文、身份验证状态、已完成步骤、失败重试历史完全一致。

\begin{importantbox}{渠道坍缩的工程含义}
渠道坍缩意味着组织模型变化：从“渠道团队优化局部指标”转向“Agent 团队优化端到端结果”。如果组织还按渠道切割，系统很难做到单一决策平面和统一记忆。
\end{importantbox}

\begin{warningbox}{迁移中的 KPI 陷阱}
很多企业会在初期把语音与在线聊天分别考核，导致 Agent 团队为了局部指标做冲突优化。结果是跨渠道一致性下降，用户重复解释问题，整体满意度反而下滑。
\end{warningbox}

\subsection{本章小结}

“单 Agent、多入口”是产品形态，不是架构捷径。要兑现这个形态，团队必须同步改造策略层、状态层、组织边界和评测指标。

